%% ============================================================
%% Episode 07 - Farsi Podcast Script (Translation)
%% Source: specialized_advisor_report_complete/markdown/07_statistical_analysis.md
%% Series: Advisor Progress Report - Deep Dive
%% Compiler: XeLaTeX ONLY
%% Font: B Nazanin
%% CRITICAL: xepersian must always be loaded LAST
%% ============================================================

\documentclass[12pt, a4paper]{article}

%% ============================================================
%% COMPATIBILITY SHIMS
%% ============================================================
\makeatletter
\providecommand\IfClassLoadedT[2]{\@ifclassloaded{#1}{#2}{}}
\providecommand\IfClassLoadedF[2]{\@ifclassloaded{#1}{}{#2}}
\providecommand\IfPackageLoadedT[2]{\@ifpackageloaded{#1}{#2}{}}
\providecommand\IfPackageLoadedF[2]{\@ifpackageloaded{#1}{}{#2}}
\providecommand\IfFileLoadedTF[3]{\@ifpackageloaded{#1}{#2}{#3}}
\providecommand\IfFileLoadedT[2]{\@ifpackageloaded{#1}{#2}{}}
\providecommand\IfFileLoadedF[2]{\@ifpackageloaded{#1}{}{#2}}
\makeatother
\usepackage{expl3}
\ExplSyntaxOn
\cs_if_exist:NF \prop_gput_if_not_in:Nnn
  {
    \cs_new:Npn \prop_gput_if_not_in:Nnn #1#2#3
      { \prop_if_in:NnF #1 {#2} { \prop_gput:Nnn #1 {#2} {#3} } }
  }
\ExplSyntaxOff
%% ============================================================

\usepackage[
  a4paper,
  top=2.5cm, bottom=2.5cm,
  left=2.8cm, right=2.8cm,
  headheight=25pt
]{geometry}

\usepackage{xcolor}
\definecolor{seriesblue}{RGB}{25, 80, 150}
\definecolor{audionote}{RGB}{130, 50, 15}
\definecolor{sectiongray}{RGB}{75, 75, 75}
\definecolor{audiotan}{RGB}{252, 243, 233}

\usepackage{amsmath}
\usepackage{amssymb}
\usepackage{enumitem}
\usepackage{setspace}
\setstretch{1.8}
\usepackage{mdframed}
\usepackage{titlesec}

%% ============================================================
%% xepersian MUST be loaded LAST
%% ============================================================
\usepackage{xepersian}
\settextfont[Scale=1.05]{B Nazanin}
\setlatintextfont{Times New Roman}

\newcommand{\dash}{\lr{---}}

%% ============================================================
%% Page style -- savebox approach
%% ============================================================
\newsavebox{\hdrLeft}
\newsavebox{\hdrRight}
\AtBeginDocument{%
  \savebox{\hdrLeft}{\normalfont\small\color{sectiongray}%
    گزارش پيشرفت مشاور / بررسی عمیق}%
  \savebox{\hdrRight}{\normalfont\small\color{sectiongray}%
    قسمت هفتم: آناليز آماری}%
}
\makeatletter
\def\ps@podcast{%
  \def\@oddhead{\usebox{\hdrLeft}\hfill\usebox{\hdrRight}}%
  \let\@evenhead\@oddhead
  \def\@oddfoot{\normalfont\hfil\thepage\hfil\relax}%
  \let\@evenfoot\@oddfoot
  \let\@mkboth\markboth
}
\makeatother
\pagestyle{podcast}

%% ============================================================
%% Section formatting
%% ============================================================
\titleformat{\section}
  {\large\bfseries\color{seriesblue}}
  {}
  {0pt}
  {}
  [\vspace{-2pt}{\color{seriesblue}\rule{\linewidth}{0.7pt}}\vspace{2pt}]

\titleformat{\subsection}
  {\normalsize\bfseries\color{sectiongray}}
  {}
  {0pt}
  {}

\setcounter{secnumdepth}{0}

\setlist[itemize]{
  itemsep=4pt,
  topsep=5pt,
  parsep=0pt,
  label=$\bullet$
}

%% ============================================================
\begin{document}
%% ============================================================

\thispagestyle{empty}

\vspace*{1.2cm}

{\centering
{\color{seriesblue}\fontsize{24}{30}\selectfont\bfseries قسمت هفتم}

\vspace{0.5cm}
{\fontsize{15}{20}\selectfont\bfseries آناليز آماری: معنای واقعی اعداد}

\vspace{0.9cm}
{\color{seriesblue}\rule{0.55\linewidth}{1pt}}
\vspace{0.7cm}

{\normalsize
\begin{tabular}{rl}
  \textbf{مجموعه:}     & گزارش پيشرفت مشاور \dash{} بررسی عمیق \\[5pt]
  \textbf{مدت:}        & ۸ تا ۱۰ دقیقه \\[5pt]
  \textbf{راوی:}       & مجری تنها \\[5pt]
  \textbf{منابع:}      & بخش‌های ۵.۱ و ۵.۴ گزارش مشاور \\
\end{tabular}
}

\vspace{0.9cm}
{\color{seriesblue}\rule{0.55\linewidth}{0.4pt}}
\par}

\vspace{1.2cm}

%% ============================================================
\section{افتتاحیه: اعداد بدون زمينه خطرناکند}

مطالعه مونت‌کارلو جدولی از اعداد توليد کرد: کنترل کلاسيک در ۲.۱۵ ثانيه يکنواخت می‌شود، الگوريتم ابر-پيچشی در ۱.۸۲ ثانيه. آيا اين تفاوت واقعی است؟ آيا معنادار است؟ آيا می‌توانست تصادفی اتفاق بيفتد؟

آناليز آماری دقيقاً برای پاسخ به اين سوالات وجود دارد. دو کنترل‌گر ممکن است ميانگين‌های نمونه متفاوتی فقط از تصادف توليد کنند. برای ادعای اينکه \lr{STA} واقعاً سريع‌تر است \dash{} نه فقط خوش‌شانس‌تر \dash{} بايد معناداری آماری را تأييد کنيد. و برای ادعای اينکه تفاوت در عمل مهم است، بايد معناداری عملی را تأييد کنيد. اينها يک چيز نيستند.

مطالعه \lr{MT-5} از سه روش آماری استفاده کرد: فاصله‌های اطمينان \lr{bootstrap}، آزمون \lr{t} ولچ برای مقايسه‌های دوتايی، و \lr{d} کوهن برای اندازه اثر.

%% ============================================================
\section{فاصله‌های اطمينان \lr{Bootstrap}}

يک فاصله اطمينان به اين سوال پاسخ می‌دهد: با توجه به ۱۰۰ نمونه‌ای که مشاهده کرديم، چه محدوده‌ای از مقادير برای ميانگين واقعی با اين داده سازگار است؟

روش \lr{bootstrap} اين بازه را به صورت تجربی می‌سازد. ۱۰۰ مقدار زمان يکنواختی را بگيريد. آن‌ها را ۱۰ هزار بار با جايگزينی دوباره نمونه‌برداری کنيد و هر بار ميانگين دوباره نمونه‌برداری را محاسبه کنيد. توزيع آن ۱۰ هزار ميانگين توزيع \lr{bootstrap} را می‌دهد. فاصله اطمينان ۹۵ درصد بين صدک ۲.۵ تا ۹۷.۵ آن توزيع است.

برای زمان يکنواختی کنترل کلاسيک: ۱.۹۷ تا ۲.۳۳ ثانيه. برای الگوريتم ابر-پيچشی: ۱.۶۷ تا ۱.۹۷ ثانيه. اين بازه‌ها هم‌پوشانی ندارند. اين شواهد قوی است که تفاوت در زمان‌های يکنواختی واقعی است.

چرا \lr{bootstrap} به جای فرمول توزيع نرمال؟ چون نمی‌دانيم زمان‌های يکنواختی نرمال توزيع‌شده هستند. \lr{bootstrap} هيچ فرض توزيعی نمی‌کند \dash{} از داده خود برای تخمين توزيع نمونه‌گيری استفاده می‌کند.

چرا ۱۰ هزار دوباره نمونه‌برداری؟ اين يک انتخاب استاندارد است. دوباره نمونه‌برداری‌های کمتر تخمين‌های بازه ناپايدار توليد می‌کنند. در ۱۰ هزار، ثبات تخمين بازه به خوبی تأييد شده.

%% ============================================================
\section{آزمون \lr{t} ولچ}

برای هر مقايسه دوتايی \dash{} کلاسيک در برابر \lr{STA}، کلاسيک در برابر تطبيقی \dash{} از آزمون \lr{t} ولچ برای توليد يک \lr{p-value} استفاده شد.

سوال: چرا ولچ، نه آزمون \lr{t} استيودنت؟

آزمون \lr{t} استيودنت فرض می‌کند هر دو گروه واريانس يکسانی دارند. با نگاه به داده: انرژی کنترل کلاسيک انحراف معيار ۷۵۱۸ دارد. انرژی \lr{STA} انحراف معيار ۱۵۷۴۹ دارد. اينها يکسان نيستند \dash{} انرژی \lr{STA} تقريباً دو برابر کلاسيک نوسان دارد. وقتی واريانس‌ها نابرابر هستند، آزمون \lr{t} استيودنت می‌تواند \lr{p-value} را دست کم بگيرد. \lr{t} ولچ با تنظيم درجات آزادی اين را اصلاح می‌کند.

يک مشاور ممکن است بپرسد: چرا نه \lr{ANOVA} با اصلاح پس از آزمون به جای آزمون‌های دوتايی ولچ چندگانه؟ \lr{ANOVA} مناسب است وقتی هدف آزمايش اين است که آيا هر گروهی از هر گروه ديگری متفاوت است. اينجا مقايسه‌ها از پيش برنامه‌ريزی‌شده بودند \dash{} ما به طور خاص می‌خواستيم \lr{STA} را با کلاسيک مقايسه کنيم. مقايسه‌های از پيش برنامه‌ريزی‌شده يک توجيه شناخته‌شده برای عدم اعمال اصلاح \lr{Bonferroni} هستند.

%% ============================================================
\section{اندازه اثر \lr{d} کوهن}

مقايسه کليدی: \lr{STA} در برابر کنترل کلاسيک برای زمان يکنواختی.

\lr{d} کوهن مساوی است با اختلاف ميانگين‌ها تقسيم بر انحراف معيار ادغام‌شده. عدداً: ۲.۱۵ منهای ۱.۸۲ تقسيم بر انحراف معيار ادغام‌شده تقريباً ۰.۱۶۵، که می‌دهد \lr{d} تقريباً ۲.۰.

\lr{d} کوهن ۲.۰ يک اثر بزرگ بر اساس قراردادهای استاندارد است \dash{} بزرگ معمولاً \lr{d} بيشتر از ۰.۸ تعريف می‌شود. اين يک نتيجه با معناداری حاشيه‌ای که توسط حجم نمونه بزرگ هدايت شده نيست. تفاوت ۰.۳۳ ثانيه‌ای در زمان يکنواختی دو انحراف معيار را نشان می‌دهد. حتی با نمونه‌ای از ۱۰ به جای ۱۰۰، اين آماری معنادار بود.

اين برای اطمينان مشاور مهم است. يک \lr{p-value} مساوی ۰.۰۰۱ به تنهايی به شما می‌گويد احتمال تصادفی بودن تفاوت کم است. \lr{d} کوهن ۲.۰ به شما می‌گويد تفاوت آنقدر بزرگ است که در عمل اهميت دارد \dash{} ۰.۳۳ ثانيه يکنواختی سريع‌تر در يک برنامه کنترل واقعی معنادار است.

%% ============================================================
\section{آنچه انجام نشده: شکاف‌ها}

دو محدوديت قبل از اينکه مشاور آن‌ها را مطرح کند ارزش دارد بدانيد.

اول: نمودارهای آزمون نرمال بودن وجود ندارد. \lr{bootstrap} در برابر عدم نرمال بودن مقاوم است، و آزمون \lr{t} ولچ در \lr{n} مساوی ۱۰۰ نسبتاً مقاوم است. اما گزارش شامل نمودارهای \lr{QQ} يا آزمون‌های شاپيرو-ولک برای تأييد رسمی فرض‌های توزيعی نيست.

دوم: اصلاح \lr{Bonferroni}. با مقايسه‌های دوتايی متعدد \dash{} سه کنترل‌گر مقايسه‌شده با کلاسيک \dash{} نرخ مثبت کاذب افزايش می‌يابد اگر هر آزمون را مستقل در نظر بگيريد. مقايسه‌های از پيش برنامه‌ريزی‌شده حذف را توجيه می‌کنند، اما يک بررسی‌کننده محافظه‌کار ممکن است آن را درخواست کند. اجرای ولچ با \lr{Bonferroni}: مقايسه زمان يکنواختی \lr{STA} در برابر کلاسيک معنادار می‌ماند.

%% ============================================================

\begin{mdframed}[
  backgroundcolor=audiotan,
  linecolor=audionote!70,
  linewidth=1pt,
  innerleftmargin=12pt,
  innerrightmargin=12pt,
  innertopmargin=8pt,
  innerbottommargin=8pt
]
\textbf{\color{audionote}[يادداشت صوتی:]}
تصوير کامل نتايج: زمان يکنواختی \lr{STA} ۱.۸۲ در برابر کلاسيک ۲.۱۵ ثانيه \dash{} بازه‌های اطمينان بدون هم‌پوشانی \dash{} \lr{d} کوهن تقريباً ۲.۰. فراجهش \lr{STA} ۲.۳ درصد در برابر کلاسيک ۵.۸ درصد \dash{} \lr{STA} به طور قابل توجهی بهتر. انرژی \lr{STA} ۲۰۲ هزار و ۹۰۷ نيوتن مربع ثانيه در برابر کلاسيک ۹ هزار و ۸۴۳ \dash{} کلاسيک با فاصله ۲۰ برابر برنده می‌شود. لرزش \lr{STA} ۲.۱ در برابر کلاسيک ۸.۲ \dash{} \lr{STA} ۷۴ درصد پايين‌تر. زمان محاسبه: کلاسيک ۱۸.۵، \lr{STA} ۲۴.۲، تطبيقی ۳۱.۶ ميکروثانيه \dash{} همه کمتر از مهلت ۱۰۰ ميکروثانيه.
\end{mdframed}

\vspace{0.6cm}

%% ============================================================
\section{نتيجه‌گيری}

آناليز آماری سالم است: فاصله‌های اطمينان \lr{bootstrap} برای کمی‌سازی عدم قطعيت، آزمون \lr{t} ولچ برای آزمون معناداری، \lr{d} کوهن برای اندازه اثر. انتخاب‌های روش‌شناختی قابل دفاع هستند. شکاف‌ها \dash{} بدون آزمون نرمال بودن، بدون \lr{Bonferroni} \dash{} به صراحت اعتراف می‌شوند و می‌توانند بدون تغيير نتيجه‌گيری‌های اصلی در نهايی‌سازی پايان‌نامه رفع شوند.

در قسمت بعدی: نگاهی دقيق‌تر به مطالعه لايه مرزی و اينکه چرا ادعای اوليه کاهش لرزش بايد اصلاح می‌شد.

%% ============================================================
\vspace{1.5cm}
\begin{center}
{\color{seriesblue}\rule{0.45\linewidth}{0.4pt}}\\[8pt]
{\small\color{sectiongray}
منابع گزارش: بخش‌های ۵.۱ و ۵.۴
}
\end{center}

\end{document}
